\section{方差平衡规模 $m_\times$、边际收益与非对称性}
\label{sec:part05}

\subsection{方差平衡规模的定义与物理几何意义}

在方差闭式解 $V(m) = v_R + A/m$ 中，可约的抽样随机性 $A/m$ 随规模线性衰减，而不可约的环境方差 $v_R$ 保持恒定。定义二者相等的交点规模：

\begin{definition}[方差平衡规模 $\mx$]
定义使得可约抽样方差与不可约环境方差严格相等的特征观测规模为：
\begin{equation}
  \frac{A}{\mx} = v_R \quad \iff \quad \mx \equiv \frac{A}{v_R} = \frac{\Rbar + \frac{\rho(\Rbar^2 + v_R)}{k}}{v_R}.
  \label{eq:mx-def}
\end{equation}
\end{definition}

借助平衡规模 $\mx$，总方差可紧凑表征为无量纲规模比值 $m / \mx$ 的函数：
\begin{equation}
  V(m) = v_R \left( 1 + \frac{\mx}{m} \right).
  \label{eq:var-compact}
\end{equation}

\begin{insightbox}{审慎定位：平衡规模不是临床或预警的“充分阈值”}
本手册与论文在此做出严格强调：方差平衡规模 $\mx$ 仅表示抽样项与环境项相等的方差结构转折基准（Balance point），\textbf{绝非}满足特定估计精度或预警检测功效要求的最低实用规模（Minimum operational threshold）。在不同公共卫生任务中：
\begin{itemize}
  \item 若目标为以 80\% 功效检测 15\% 的传播增长，所需规模 $m \approx 527$，约为 $\mx$ 的 10 倍（见第~\ref{sec:part11} 节）；
  \item 若目标为压制跨期宏观趋势估计误差，当规模达到 $\mx$ 时，继续扩样的边际收益已大幅钝化，资源应当向增加期数或空间池化转移。
\end{itemize}
\end{insightbox}

\subsection{边际收益递减的四大数学定理证明}

\begin{derivationbox}{推导 5.1：方差份额、翻倍收益与弹性分析}
基于式~\eqref{eq:var-compact}，我们可以严格导出方差曲线的四个关键几何定理：

\begin{enumerate}
  \item \textbf{抽样方差占比定理：}
  抽样随机性在总方差中所占的百分比为：
  \begin{equation}
    S(m) \equiv \frac{A/m}{V(m)} = \frac{A/m}{v_R + A/m} = \frac{\frac{A/v_R}{m}}{1 + \frac{A/v_R}{m}} = \frac{\mx/m}{1 + \mx/m} = \frac{1}{1 + \frac{m}{\mx}}.
    \label{eq:share-proof}
  \end{equation}
  当 $m = \mx$ 时，$S(\mx) = \frac{1}{1 + 1} = \frac{1}{2} = 50\%$。两类方差严格均分总误差。

  \item \textbf{规模翻倍收益衰减定理（Doubling reduction）：}
  将单期监测规模从 $m$ 翻倍扩大到 $2m$ 时，总方差相对下降的比例为：
  \begin{align}
    \Delta_2(m) &\equiv \frac{V(m) - V(2m)}{V(m)} = \frac{\left(v_R + \frac{A}{m}\right) - \left(v_R + \frac{A}{2m}\right)}{v_R + \frac{A}{m}} \nonumber \\
    &= \frac{\frac{A}{2m}}{v_R + \frac{A}{m}} = \frac{1}{2} \cdot \frac{A/m}{v_R + A/m} = \frac{1}{2} S(m) = \frac{1}{2\left(1 + \frac{m}{\mx}\right)}.
    \label{eq:doubling-proof}
  \end{align}
  考察三个典型区域：
  \begin{itemize}
    \item 当 $m \to 0$（小样本极端）时，$\Delta_2 \to 50\%$，符合经典泊松抽样方差减半直觉；
    \item 当 $m = \mx$（平衡点）时，$\Delta_2 = \frac{1}{4} = 25.0\%$；
    \item 当 $m = 4\mx$（大样本饱和）时，$\Delta_2 = \frac{1}{2(1+4)} = 10.0\%$，翻倍投入仅换来不到一成的方差压缩。
  \end{itemize}

  \item \textbf{规模弹性（Elasticity）定理：}
  总方差关于监测规模 $m$ 的点弹性为：
  \begin{equation}
    \epsilon_m \equiv \frac{\mathrm{d}\ln V(m)}{\mathrm{d}\ln m} = \frac{m}{V(m)} \frac{\mathrm{d}V(m)}{\mathrm{d}m} = \frac{m}{v_R + A/m} \left( -\frac{A}{m^2} \right) = -\frac{A/m}{v_R + A/m} = -\frac{\mx}{m + \mx}.
  \end{equation}
  在平衡点 $m = \mx$ 处，规模弹性严格等于 $-\frac{1}{2}$。

  \item \textbf{处处严格凸性（Strict convexity）：}
  求 $V(m)$ 关于 $m$ 的一阶与二阶导数：
  \begin{equation}
    V'(m) = -\frac{A}{m^2} < 0, \qquad V''(m) = \frac{2A}{m^3} > 0, \quad \forall m > 0.
  \end{equation}
  一阶导数处处严格为负，二阶导数在正实轴上处处严格为正。函数没有拐点，曲线呈现处处向原点凸起的双曲线单调下降特征。
\end{enumerate}
\end{derivationbox}

\subsection{覆盖面与报告率的非对称边际效应（Remark 1 证明）}

预期报告规模由两个途径共同构成：$m = \rho n$。究竟应该追求让既定人群的报告更全面（提高报告概率 $\rho$），还是覆盖更多的人口群体（扩大祖代样本 $n$）？

\begin{theorembox}{注 1 的完全证明（覆盖面与报告率的边际非对称性）}
在假设 H1--H4 下：
\begin{enumerate}[label=(\roman*)]
  \item 若固定祖代规模 $n$ 而令报告率 $\rho \to 1$，方差趋于大于 $v_R$ 的严格常数；
  \item 若固定报告率 $\rho$ 而令覆盖规模 $n \to \infty$，方差趋于理论底板 $v_R$；
  \item 按相同百分比增加 $n$ 所取得的方差降幅，恒等于增加 $\rho$ 所取得降幅的 $A/\Rbar$ 倍：
  \begin{equation}
    \frac{-\frac{\partial V}{\partial \ln n}}{-\frac{\partial V}{\partial \ln \rho}} = \frac{A}{\Rbar} = 1 + \frac{\rho(\Rbar^2 + v_R)}{k\Rbar} > 1.
    \label{eq:asym-ratio}
  \end{equation}
\end{enumerate}
\end{theorembox}

\begin{proof}
将方差式~\eqref{eq:var-closed} 展开为 $n$ 和 $\rho$ 的显式二元函数：
\begin{equation}
  V(n, \rho) = v_R + \frac{\Rbar}{\rho n} + \frac{\Rbar^2 + v_R}{kn}.
\end{equation}
\begin{enumerate}
  \item 固定 $n$ 令 $\rho \to 1$：
  \begin{equation}
    \lim_{\rho \to 1} V(n, \rho) = v_R + \frac{\Rbar}{n} + \frac{\Rbar^2 + v_R}{kn} = v_R + \frac{1}{n}\left(\Rbar + \frac{\Rbar^2+v_R}{k}\right) > v_R.
  \end{equation}
  因超传播项不随 $\rho$ 衰减，存在由有限祖代规模 $n$ 锁定的残余人口随机性。
  \item 固定 $\rho$ 令 $n \to \infty$：
  \begin{equation}
    \lim_{n \to \infty} V(n, \rho) = v_R + 0 + 0 = v_R.
  \end{equation}
  \item 计算对数偏导数（即方差对比例变化的敏感度）：
  \begin{align}
    -\frac{\partial V}{\partial \ln n} &= -n \frac{\partial V}{\partial n} = -n \left[ -\frac{\Rbar}{\rho n^2} - \frac{\Rbar^2 + v_R}{kn^2} \right] = \frac{\Rbar}{\rho n} + \frac{\Rbar^2 + v_R}{kn} = \frac{A}{m}, \\
    -\frac{\partial V}{\partial \ln \rho} &= -\rho \frac{\partial V}{\partial \rho} = -\rho \left[ -\frac{\Rbar}{\rho^2 n} \right] = \frac{\Rbar}{\rho n} = \frac{\Rbar}{m}.
  \end{align}
  两者之比为：
  \begin{equation}
    \frac{-\partial V / \partial \ln n}{-\partial V / \partial \ln \rho} = \frac{A/m}{\Rbar/m} = \frac{A}{\Rbar} = \frac{\Rbar + \frac{\rho(\Rbar^2+v_R)}{k}}{\Rbar} = 1 + \frac{\rho(\Rbar^2 + v_R)}{k\Rbar}.
  \end{equation}
  由于 $\rho > 0, k > 0, \Rbar > 0, v_R \ge 0$，第二项严格为正，故比值恒大于 1。证毕。
\end{enumerate}
\end{proof}
